\documentclass[10pt,a4paper]{amsart}
\usepackage{amsmath,amssymb,amsthm,graphicx}
\usepackage{verbatim}
\newenvironment{reference}{\comment}{\endcomment}
%\newenvironment{reference}{}{}
\newenvironment{slogan}{\comment}{\endcomment}
\newenvironment{history}{\comment}{\endcomment}
% For commutative diagrams we use Xy-pic
\usepackage[all]{xy}

% We use 2cell for 2-commutative diagrams.
\xyoption{2cell}
% Theorem environments.
%
\theoremstyle{plain}
\newtheorem{theorem}[subsection]{Theorem}
\newtheorem{proposition}[subsection]{Proposition}
\newtheorem{lemma}[subsection]{Lemma}

\theoremstyle{definition}
\newtheorem{definition}[subsection]{Definition}
\newtheorem{example}[subsection]{Example}
\newtheorem{exercise}[subsection]{Exercise}
\newtheorem{situation}[subsection]{Situation}

\theoremstyle{remark}
\newtheorem{remark}[subsection]{Remark}
\newtheorem{remarks}[subsection]{Remarks}

\numberwithin{equation}{subsection}

% Macros
%
\def\lim{\mathop{\mathrm{lim}}\nolimits}
\def\colim{\mathop{\mathrm{colim}}\nolimits}
\def\Spec{\mathop{\mathrm{Spec}}}
\def\Hom{\mathop{\mathrm{Hom}}\nolimits}
\def\Ext{\mathop{\mathrm{Ext}}\nolimits}
\def\SheafHom{\mathop{\mathcal{H}\!\mathit{om}}\nolimits}
\def\SheafExt{\mathop{\mathcal{E}\!\mathit{xt}}\nolimits}
\def\Sch{\mathit{Sch}}
\def\Mor{\mathop{\mathrm{Mor}}\nolimits}
\def\Ob{\mathop{\mathrm{Ob}}\nolimits}
\def\Sh{\mathop{\mathit{Sh}}\nolimits}
\def\NL{\mathop{N\!L}\nolimits}
\def\CH{\mathop{\mathrm{CH}}\nolimits}
\def\proetale{{pro\text{-}\acute{e}tale}}
\def\etale{{\acute{e}tale}}
\def\QCoh{\mathit{QCoh}}
\def\Ker{\mathop{\mathrm{Ker}}}
\def\Im{\mathop{\mathrm{Im}}}
\def\Coker{\mathop{\mathrm{Coker}}}
\def\Coim{\mathop{\mathrm{Coim}}}

% Boxtimes
%
\DeclareMathSymbol{\boxtimes}{\mathbin}{AMSa}{"02}

%
% Macros for moduli stacks/spaces
%
\def\QCohstack{\mathcal{QC}\!\mathit{oh}}
\def\Cohstack{\mathcal{C}\!\mathit{oh}}
\def\Spacesstack{\mathcal{S}\!\mathit{paces}}
\def\Quotfunctor{\mathrm{Quot}}
\def\Hilbfunctor{\mathrm{Hilb}}
\def\Curvesstack{\mathcal{C}\!\mathit{urves}}
\def\Polarizedstack{\mathcal{P}\!\mathit{olarized}}
\def\Complexesstack{\mathcal{C}\!\mathit{omplexes}}
% \Pic is the operator that assigns to X its picard group, usage \Pic(X)
% \Picardstack_{X/B} denotes the Picard stack of X over B
% \Picardfunctor_{X/B} denotes the Picard functor of X over B
\def\Pic{\mathop{\mathrm{Pic}}\nolimits}
\def\Picardstack{\mathcal{P}\!\mathit{ic}}
\def\Picardfunctor{\mathrm{Pic}}
\def\Deformationcategory{\mathcal{D}\!\mathit{ef}}
\title{The scheme theoretic image is insensitive to a surjective flat cover of finite presentation}
\author{The Stacks Project Authors}
\date{stacks-project:stacks-morphisms.tex:local-snapshot-20261012; GNU FDL 1.2 or later}
\begin{document}
\maketitle
\noindent\textit{Editorial note.} All delivered mathematical text is the author's own, copied from the pinned Stacks Project chapter source sources/stacks-algebraic-stacks-morphisms/pinned-original-source.tex; no statement or proof step was rewritten, reordered or removed.\par\medskip
\noindent\textit{Editorial note.} Preamble provenance: the Stacks Project keeps one single shared preamble for every chapter; this card takes it directly from the staged sources/stacks-algebraic-stacks-morphisms/preamble.tex as three verbatim spans (lines 8-12, 14-18 and 154-217), so no derived file was written and nothing was added or redefined. The wrapper class and the base packages are supplied by the assembler.\par\medskip
\noindent\textit{Editorial note.} Bibliography provenance: neither the counted statement nor its proof contains a citation (this lane rescans the whole span before writing the card), so the card carries no bibliography and no entry was reconstructed.\par\medskip
\noindent\textit{Editorial note.} Selection provenance: this card was opened under the file-based control contract (CONTROL.md), with the mechanical pre-card gate run on its statement and proof interval before the card was written.\par\medskip
\noindent\textit{Allowed standard prerequisites (editorial).}\begin{itemize}
\item The scheme theoretic image of a morphism of algebraic stacks, characterised by the closed substacks through which the morphism factors. The author's section on scheme theoretic images, in which the counted statement occurs. Supplies the universal property against which existence and equality are tested.
\item The fibre product of a morphism along a surjective flat morphism of finite presentation is again surjective flat of finite presentation, so that it admits an fppf covering by a scheme. Used inside the counted proof (source lines 8768-8776). Produces the covering of the test scheme by which the object is tested for membership in a closed substack.
\end{itemize}\medskip
\begin{lemma}
\label{lemma-cover-upstairs}
Let $f : \mathcal{X} \to \mathcal{Y}$ be a morphism of algebraic stacks.
Let $g : \mathcal{W} \to \mathcal{X}$ be a morphism of algebraic stacks
which is surjective, flat, and locally of finite presentation.
Then the scheme theoretic image of $f$ exists if and only if the
scheme theoretic image of $f \circ g$ exists and if so then these
scheme theoretic images are the same.
\end{lemma}

\begin{proof}
Assume $\mathcal{Z} \subset \mathcal{Y}$
is a closed substack and $f \circ g$ factors through $\mathcal{Z}$.
To prove the lemma it suffices to show
that $f$ factors through $\mathcal{Z}$.
Consider a scheme $T$ and a morphism $T \to \mathcal{X}$
given by an object $x$ of the fibre category of $\mathcal{X}$ over $T$.
We will show that $f(x)$ is in fact in the fibre category of $\mathcal{Z}$
over $T$. Namely, the projection $T \times_\mathcal{X} \mathcal{W} \to T$
is a surjective, flat, locally finitely presented morphism.
Hence there is an fppf covering $\{T_i \to T\}$ such that
$T_i \to T$ factors through $T \times_\mathcal{X} \mathcal{W} \to T$
for all $i$. Then $T_i \to \mathcal{X}$ factors through $\mathcal{W}$
and hence $T_i \to \mathcal{Y}$ factors through $\mathcal{Z}$.
Thus $f(x)|_{T_i}$ is an object of $\mathcal{Z}$.
Since $\mathcal{Z}$ is a strictly full substack, we conclude
that $f(x)$ is an object of $\mathcal{Z}$ as desired.
\end{proof}

\end{document}
